\chapter[Estatística I]{Estatística desde o começo: posição, dispersão e boxplot}

Imagine que a coordenação de um curso recebeu uma planilha com os tempos, em minutos, que oito estudantes levaram para concluir uma atividade de análise de dados: $18$, $20$, $20$, $22$, $25$, $25$, $30$ e $56$. A primeira reação pode ser procurar um botão que produza um painel. Entretanto, antes do gráfico existe uma sequência de decisões: esses oito estudantes representam quem? O tempo foi medido da mesma maneira para todos? O valor $56$ é erro, caso raro ou observação legítima? Qual número resume o grupo sem esconder sua diversidade? A estatística nasceu justamente da necessidade humana de transformar muitas observações em descrições comparáveis sem perder de vista aquilo que foi observado.

Neste capítulo, começaremos do início absoluto. Primeiro distinguiremos população, amostra, unidade observada, variável e frequência. Depois construiremos, uma por vez, as medidas de posição: média simples, média ponderada, mediana, moda, quartis e percentis. Em seguida, perceberemos por que um centro não basta e chegaremos à amplitude, à variância populacional, à variância amostral, ao desvio-padrão e ao intervalo interquartil. O boxplot fechará a história ao reunir posição e dispersão em uma representação compacta. Não usaremos código: todas as contas serão abertas e acompanhadas linha a linha.

O fio condutor será a turma fictícia da Trilha UNIFACOL. Trabalharemos com conjuntos pequenos porque eles permitem enxergar cada decisão e conferir cada operação. Isso não torna a matemática artificial. Uma fórmula aplicada a milhões de registros executa o mesmo raciocínio que faremos com cinco ou oito valores; a diferença é apenas o volume. Aprender com poucos dados cria uma referência contra a qual qualquer ferramenta poderá ser verificada posteriormente.

O encontro foi dimensionado para três horas. No primeiro bloco, construímos o vocabulário e as tabelas de frequência; no segundo, localizamos o centro com média, mediana e moda; depois da chamada das 20h30, medimos posição relativa e dispersão; no fechamento, montamos e interpretamos um boxplot. As contas resolvidas fazem parte da explicação do livro, não constituem uma lista de exercícios. A intenção é que o professor possa refazer cada linha com a turma enquanto os estudantes explicam o significado, e não somente repetem uma regra.

Ao final, você deverá conseguir escolher uma medida de acordo com a pergunta, explicar cada símbolo das fórmulas e reconstruir os cálculos sem depender de uma função pronta. Também deverá distinguir variância populacional e amostral, justificar o divisor $n-1$ sem transformar a justificativa em mistério e interpretar o que o desvio-padrão e o IQR dizem na unidade do problema. Mais importante: deverá perceber que estatística descritiva não é uma coleção de botões. É uma linguagem para reduzir dados sem apagar contexto.

\begin{SummaryBox}
Então, aluno, veja só o caminho desta aula: começaremos perguntando quem e o que foi observado, organizaremos as ocorrências e somente depois resumiremos posição e dispersão. Cada fórmula aparecerá porque uma pergunta real exigiu uma resposta mais precisa.
\end{SummaryBox}

\section{Antes de calcular: população, amostra, variáveis e frequências}

Muito antes de existirem computadores, governos, comerciantes e pesquisadores precisavam contar pessoas, terras, colheitas e impostos. A palavra \textit{estatística} está historicamente ligada ao Estado e à produção de informações sobre uma coletividade. O problema essencial permanece atual: raramente queremos apenas saber os valores escritos em uma tabela; queremos dizer algo sobre um conjunto de pessoas, objetos ou eventos. Chamamos de \textbf{população} o conjunto completo de unidades sobre o qual a pergunta foi formulada. Se a pergunta é ``qual foi o tempo de conclusão de todos os estudantes matriculados em AED neste semestre?'', a população inclui todos os estudantes que satisfazem esse recorte, e não apenas aqueles que responderam primeiro.

Nem sempre é possível observar toda a população. Pode ser caro, demorado, destrutivo ou simplesmente inviável. Selecionamos então uma \textbf{amostra}, isto é, uma parte da população efetivamente observada para produzir conhecimento sobre o conjunto maior. Se a turma possui quarenta estudantes e medimos oito, temos tamanho populacional $N=40$ e tamanho amostral $n=8$. As letras maiúscula e minúscula ajudam a recordar a diferença, mas não garantem qualidade. Uma amostra de oito voluntários pode representar mal a turma se apenas os mais rápidos aceitaram participar. O modo de seleção importa tanto quanto o tamanho.

Cada elemento sobre o qual registramos informações é uma \textbf{unidade de observação}. Na tabela dos tempos, uma linha pode representar uma tentativa de um estudante, não necessariamente um estudante único. Se a mesma pessoa realizou três tentativas, temos três eventos e talvez apenas uma pessoa. Antes de qualquer média, precisamos completar a frase ``cada linha representa...''. Essa frase determina o denominador. Confundir tentativa com estudante produz um resultado numericamente correto para a tabela e conceitualmente errado para a pergunta.

Uma \textbf{variável} é uma característica que pode assumir valores diferentes entre as unidades. O curso de matrícula e o turno são qualitativos nominais: indicam categorias sem ordem natural. Uma avaliação como ``insuficiente, básico, adequado, avançado'' é qualitativa ordinal: existe ordem, mas a distância entre categorias não foi medida como igual. Número de faltas é quantitativo discreto porque resulta de contagem. Tempo, massa e temperatura são quantitativos contínuos porque, em princípio, admitem valores em um intervalo, embora o instrumento imponha precisão. Classificar a variável impede operações sem sentido, como calcular a média de códigos de curso $1$, $2$ e $3$.

Um teste semântico simples ajuda quando a coluna parece numérica: ``faz sentido somar dois valores?'' Para salários de R\$ 3.000 e R\$ 5.000, a soma R\$ 8.000 representa um total; a diferença R\$ 2.000 também possui significado. Para códigos de região $1$ e $2$, o resultado $3$ não representa a junção das regiões, e a diferença não mede distância geográfica. Pergunte ainda se o zero significa ausência da quantidade ou apenas um rótulo. Matrícula, CEP e telefone permanecem identificadores mesmo armazenados como números; tempo e faltas são quantidades. Esse teste prepara tanto a escolha da média quanto as restrições do coeficiente de variação.

\begin{GuidedBox}
Aluno, preste atenção aqui: o tipo exibido pelo sistema não decide o tipo estatístico. Antes de somar, subtrair ou calcular média, pergunte se soma, diferença e zero possuem significado no fenômeno. Um código numérico continua sendo categoria quando seus algarismos apenas identificam.
\end{GuidedBox}

Para organizar observações, começamos pela \textbf{frequência absoluta}. Dez estudantes informaram, na ordem, celular, notebook, celular, tablet, notebook, celular, notebook, notebook, celular e notebook. Contamos $4$ ocorrências de celular, $5$ de notebook e $1$ de tablet. Denotando a frequência absoluta da categoria $j$ por $f_j$, temos $f_{cel}=4$, $f_{note}=5$ e $f_{tab}=1$. A soma precisa recuperar o número de observações: $4+5+1=10=n$. Essa igualdade é uma primeira conferência do raciocínio.

A frequência relativa responde qual parte do total pertence a cada categoria. Sua fórmula é
\[
 h_j=\frac{f_j}{n}.
\]
Aqui, $h_j$ é a frequência relativa da categoria $j$, $f_j$ é sua contagem e $n$ é o total observado. Para celular, $h_{cel}=4/10=0{,}40=40\%$; para notebook, $h_{note}=5/10=0{,}50=50\%$; para tablet, $h_{tab}=1/10=0{,}10=10\%$. Somando, $0{,}40+0{,}50+0{,}10=1$, ou $100\%$. A contagem mostra volume; a proporção permite comparar grupos de tamanhos diferentes.

\begin{center}
\small
\begin{tabular}{@{}lrr@{}}
\toprule
\textbf{Canal} & \textbf{Frequência absoluta $f_j$} & \textbf{Frequência relativa $h_j$} \\
\midrule
Celular & 4 & 40\% \\
Notebook & 5 & 50\% \\
Tablet & 1 & 10\% \\
\midrule
Total & 10 & 100\% \\
\bottomrule
\end{tabular}
\end{center}

Quando as categorias possuem ordem, podemos calcular frequência acumulada. Considere as avaliações ${1,2,2,3,3,3,4,4,5,5}$. Até o valor $1$, há $1$ caso; até $2$, há $1+2=3$; até $3$, $3+3=6$; até $4$, $6+2=8$; até $5$, $8+2=10$. A frequência relativa acumulada até $3$ é $6/10=60\%$. Isso significa que seis em cada dez observações estão no máximo em $3$. A palavra ``até'' é crucial: frequência acumulada não serve naturalmente a categorias nominais como cor ou curso, porque não há uma ordem legítima a acumular.

\begin{GuidedBox}
Aluno, preste atenção aqui: população é o conjunto sobre o qual queremos falar; amostra é o conjunto que observamos; unidade é aquilo que cada linha representa; variável é a característica registrada. Se essas quatro respostas estiverem erradas, uma conta perfeita produzirá uma conclusão errada.
\end{GuidedBox}

A transição para as medidas de posição nasce de uma necessidade simples. A tabela de frequências descreve quantas vezes cada valor apareceu, mas um relatório frequentemente precisa de um resumo do valor típico. Não existe um único ``centro'' universal. A média redistribui o total igualmente, a mediana localiza a posição central e a moda identifica a maior repetição. Começaremos pela média porque sua ideia de repartição equitativa deixa visível por que cada observação participa do cálculo.

\begin{SummaryBox}
Então, aluno, veja só o que você aprendeu aqui: antes da fórmula vêm o recorte, a unidade e o tipo da variável. Frequência absoluta conta ocorrências; frequência relativa divide a contagem pelo total; frequência acumulada informa quantos casos estão até determinada posição ordenável.
\end{SummaryBox}

\section{Média simples e média ponderada: repartir o total com justiça matemática}

A média aritmética surgiu da necessidade de representar várias magnitudes por um único valor comparável. Em astronomia, medições repetidas do mesmo fenômeno variavam por erros de observação; em comércio e administração, era necessário comparar preços, produções e rendimentos. A ideia é antiga e poderosa: somamos o que foi observado e imaginamos uma redistribuição igual entre as unidades. Por isso, a média pode ser entendida como ponto de equilíbrio. Ela preserva o total, embora nenhum indivíduo precise possuir exatamente o valor médio.

Para uma amostra com valores $x_1,x_2,\ldots,x_n$, a média aritmética é
\[
 \bar{x}=\frac{1}{n}\sum_{i=1}^{n}x_i
 =\frac{x_1+x_2+\cdots+x_n}{n}.
\]
O símbolo $x_i$ representa o valor observado na posição $i$; $i$ percorre as posições de $1$ até $n$; $n$ é a quantidade de observações; $\sum$ ordena que todos os valores sejam somados; e $\bar{x}$, lido ``x barra'', representa a média da amostra. Se tivéssemos a população inteira, seria comum usar $\mu$, a letra grega mi, para a média populacional: $\mu=(1/N)\sum_{i=1}^{N}x_i$.

Vamos calcular a média das horas de estudo de cinco estudantes: $x={2,4,4,5,10}$. Primeiro contamos $n=5$. Depois somamos cada observação, sem pular etapas:
\[
\sum_{i=1}^{5}x_i=2+4+4+5+10=25.
\]
Por fim, dividimos o total pelo número de estudantes:
\[
\bar{x}=\frac{25}{5}=5\text{ horas}.
\]
O resultado diz que, se as $25$ horas fossem redistribuídas igualmente, cada estudante teria $5$ horas. Não diz que a maioria estudou cinco horas; apenas um valor observado foi $5$.

A propriedade de equilíbrio pode ser conferida pelos desvios em relação à média. Subtraímos $5$ de cada valor: $2-5=-3$, $4-5=-1$, $4-5=-1$, $5-5=0$ e $10-5=5$. Somando os desvios, $-3-1-1+0+5=0$. Os valores abaixo da média equilibram os valores acima. Essa igualdade, $\sum(x_i-\bar{x})=0$, será importante quando chegarmos à variância: ela explica por que não podemos simplesmente calcular a média dos desvios assinados para medir dispersão.

A média utiliza todos os valores e, por isso, reage fortemente aos extremos. Troque $10$ por $30$: o conjunto vira $\{2,4,4,5,30\}$. A soma passa a $45$ e a média passa a $45/5=9$. Quatro dos cinco valores permanecem no máximo em $5$, mas a média sobe de $5$ para $9$. Isso não torna a média errada. Mostra que sua pergunta é ``qual é a repartição igual do total?'', e que um valor grande altera bastante esse total. Para encontrar uma posição central resistente, precisaremos da mediana.

A Figura~\ref{fig:efeito-outlier-centro} coloca lado a lado o conjunto original e o conjunto em que apenas o último valor mudou. Os pontos representam observações, a linha azul marca a média e a linha laranja tracejada marca a mediana. A figura deve ser lida junto da conta anterior: ela torna visível qual medida se desloca quando um valor extremo altera a soma.

\begin{figure}[htbp]
\centering
\begin{tikzpicture}[x=.27cm,y=1.35cm]
\draw[->] (0,.45)--(32,.45) node[right] {valor};
\foreach \x in {0,5,...,30}{\draw[gray!55] (\x,.39)--(\x,.51) node[below] {\x};}
\node[anchor=east] at (0,2) {original};
\node[anchor=east] at (0,1) {com extremo 30};
\foreach \x in {2,4,5,10}{\fill (\x,2) circle (2.4pt);}
\fill (4,2.13) circle (2.4pt);
\foreach \x in {2,4,5,30}{\fill (\x,1) circle (2.4pt);}
\fill (4,1.13) circle (2.4pt);
\draw[guidedblueframe,very thick] (5,1.65)--(5,2.35) node[above] {$\bar x=5$};
\draw[guidedblueframe,very thick] (9,.65)--(9,1.35) node[above] {$\bar x=9$};
\draw[guidedorangeframe,very thick,dashed] (4,.65)--(4,2.35);
\end{tikzpicture}
\caption{Um único extremo desloca a média de 5 para 9, enquanto a mediana permanece em 4. Fonte: elaboração didática.}
\label{fig:efeito-outlier-centro}
\end{figure}

Observe na Figura~\ref{fig:efeito-outlier-centro} que o ponto $30$ não ``quebra'' a média: ele acrescenta vinte unidades à soma em comparação com o $10$, e essas vinte unidades são repartidas entre cinco observações, deslocando a média em quatro. A mediana permanece $4$ porque a terceira posição ordenada não mudou. Se o total interessa, a mudança da média é informação; se o centro posicional interessa, a estabilidade da mediana é informação. Ocultar uma delas empobreceria a análise.

Antes da mediana, precisamos resolver outro problema: algumas observações representam quantidades diferentes. A \textbf{média ponderada} atribui um peso $w_i$ a cada valor $x_i$:
\[
 \bar{x}_w=\frac{\sum_{i=1}^{k}w_i x_i}{\sum_{i=1}^{k}w_i}.
\]
Aqui, $k$ é o número de valores ou grupos, $w_i$ é o peso do grupo $i$, $w_ix_i$ é sua contribuição ao total e $\sum w_i$ é o peso total. A média simples é um caso particular no qual todos os pesos são iguais a $1$.

Considere duas avaliações. A atividade vale peso $2$ e a prova vale peso $8$. Um estudante obteve $9$ na atividade e $6$ na prova. Não podemos calcular $(9+6)/2=7{,}5$, porque isso daria a mesma importância às duas avaliações. A conta ponderada começa pelas contribuições:
\[
2\times9=18,\qquad 8\times6=48.
\]
Somamos as contribuições, $18+48=66$, e os pesos, $2+8=10$. Então,
\[
\bar{x}_w=\frac{66}{10}=6{,}6.
\]
O resultado fica mais próximo de $6$ porque a prova concentra $80\%$ do peso.

A ponderação também corrige a chamada ``média das médias''. Uma turma A com $10$ estudantes tem média $8$; uma turma B com $30$ estudantes tem média $6$. A média simples das turmas é $(8+6)/2=7$, adequada apenas se quisermos que cada turma tenha o mesmo peso. Para descrever o estudante típico entre os quarenta, recuperamos os totais: $10\times8=80$ pontos na turma A e $30\times6=180$ na B. Logo,
\[
\bar{x}_w=\frac{80+180}{10+30}=\frac{260}{40}=6{,}5.
\]
O valor $6{,}5$ responde à população de estudantes; $7$ responde a duas turmas tratadas como duas unidades. A pergunta determina o peso.

\begin{GuidedBox}
Aluno, preste atenção aqui: peso não é um enfeite colocado depois da conta. Ele declara quanto cada valor representa. Antes de usar média ponderada, explique a origem dos pesos e confira se o denominador é a soma deles, não apenas a quantidade de linhas.
\end{GuidedBox}

Chegamos a uma escolha importante. A média é excelente quando o total e o equilíbrio importam, como em custo médio, nota ponderada ou consumo por unidade. Entretanto, quando um extremo desloca o total, pode deixar de representar a posição onde se encontra a maioria. A próxima medida abandona a lógica de repartir e passa a perguntar: depois de ordenar, qual observação ocupa o meio?

\begin{SummaryBox}
Então, aluno, veja só o que você aprendeu aqui: a média simples reparte a soma igualmente entre $n$ observações. A média ponderada reparte a soma das contribuições pela soma dos pesos. Ambas usam todos os valores; por isso, um extremo legítimo também participa integralmente do resultado.
\end{SummaryBox}

\section{Mediana e moda: o centro pela posição e pela repetição}

A mediana atende a uma necessidade diferente da média: localizar o centro de uma sequência ordenada. A ideia tornou-se especialmente útil em distribuições assimétricas, como renda, preço de imóveis e tempo de atendimento, nas quais poucos valores muito altos puxam a média. Em vez de redistribuir o total, a mediana divide as observações por posição: aproximadamente metade fica abaixo e metade acima. Ela é uma medida de ordem, não de soma.

O primeiro passo é sempre ordenar. Para $n$ ímpar, a posição central é
\[
 p=\frac{n+1}{2}.
\]
Considere ${10,2,5,4,4}$. A ordem de coleta não serve para localizar o meio. Ordenamos: ${2,4,4,5,10}$. Como $n=5$, $p=(5+1)/2=3$. O terceiro valor é $4$; portanto, a mediana, indicada por $\widetilde{x}$ ou $Q_2$, é $4$. Dois valores estão antes e dois depois da posição central.

Quando $n$ é par, não existe uma única posição central. Calculamos a média dos valores nas posições $n/2$ e $n/2+1$. Para ${2,4,5,10}$, já ordenado, $n=4$. As posições centrais são $4/2=2$ e $4/2+1=3$. Os valores correspondentes são $4$ e $5$. Assim,
\[
\widetilde{x}=\frac{4+5}{2}=\frac{9}{2}=4{,}5.
\]
A mediana não precisa ser um valor observado; $4{,}5$ é o ponto que separa as duas observações inferiores das duas superiores.

Compare agora média e mediana no conjunto ${2,4,4,5,30}$. A média já calculada é $9$. A mediana é o terceiro valor, $4$. A diferença conta uma história: o centro por posição permanece próximo do grupo principal, enquanto o centro por equilíbrio foi deslocado pelo $30$. Se o objetivo for estimar a despesa total dividida igualmente, a média continua relevante. Se a pergunta for ``qual tempo representa a experiência central de uma pessoa?'', a mediana provavelmente comunica melhor. Estatística não escolhe automaticamente; ela torna a escolha justificável.

A \textbf{moda} nasceu da necessidade de identificar o valor, categoria ou classe mais frequente. Sua definição é direta: moda é o valor associado à maior frequência absoluta. No conjunto ${2,4,4,5,30}$, as frequências são $f(2)=1$, $f(4)=2$, $f(5)=1$ e $f(30)=1$. Logo, a moda é $4$. Diferentemente da média e da mediana, ela pode ser aplicada a uma variável nominal. Podemos dizer que notebook é o canal modal sem atribuir números artificiais às categorias.

Um conjunto pode ser amodal, unimodal, bimodal ou multimodal. Em ${1,2,3,4}$, todos aparecem uma vez; não há um valor mais frequente, então o conjunto é amodal. Em ${1,1,2,3}$, a moda única é $1$. Em ${1,1,2,2,3}$, há duas modas, $1$ e $2$. Em pesquisa aplicada, duas modas podem indicar subgrupos, como estudantes que concluem uma tarefa em torno de vinte minutos e outro grupo em torno de quarenta. Resumir tudo por uma média poderia esconder essa estrutura.

Moda não significa automaticamente ``melhor'' nem ``recomendado''. Ela informa repetição. Se o tamanho de camiseta modal é M, o estoque precisa considerar a frequência de M, mas não pode ignorar P e G. Se o erro mais frequente de um sistema é ``senha inválida'', isso indica volume, não necessariamente maior impacto. Uma falha rara de segurança pode exigir prioridade superior. A medida responde ``qual ocorre mais?'', e a decisão acrescenta custo, risco e finalidade.

Podemos reunir as três medidas no conjunto de tempos ${18,20,20,22,25,25,30,56}$. A soma é $216$, então a média é $216/8=27$. As posições centrais são a quarta e a quinta, $22$ e $25$, logo a mediana é $(22+25)/2=23{,}5$. As maiores frequências são as de $20$ e $25$, ambas iguais a $2$, portanto a distribuição é bimodal. O mesmo conjunto possui três centros legítimos porque cada um responde a uma pergunta diferente.

\begin{center}
\small
\begin{tabularx}{0.94\textwidth}{@{}p{2.6cm}YY@{}}
\toprule
\textbf{Medida} & \textbf{Pergunta que responde} & \textbf{Resultado nos tempos} \\
\midrule
Média & qual seria a repartição igual do total? & $27$ minutos \\
Mediana & qual é o centro das posições ordenadas? & $23{,}5$ minutos \\
Moda & quais valores mais se repetem? & $20$ e $25$ minutos \\
\bottomrule
\end{tabularx}
\end{center}

A mediana prepara a próxima ideia. Se conseguimos dividir os dados em duas metades, também podemos dividi-los em quatro partes ou em cem partes. Quartis e percentis generalizam a posição relativa e permitem dizer não apenas onde está o centro, mas onde uma observação se encontra em relação às demais.

\begin{SummaryBox}
Então, aluno, veja só o que você aprendeu aqui: mediana exige ordenar e localiza o meio; moda identifica a maior frequência. Média, mediana e moda não competem para ver qual é sempre melhor. Cada uma preserva uma característica diferente dos dados.
\end{SummaryBox}

\section{Quartis e percentis: localizar uma observação na distribuição}

Quartis e percentis respondem a uma necessidade administrativa e científica antiga: localizar uma observação dentro de uma distribuição. Em educação, saúde, economia e controle de qualidade, não basta saber o valor central; queremos reconhecer faixas e posições relativas. Um percentil não diz apenas ``quanto alguém obteve'', mas ``que parcela das observações está até esta posição''. Essa linguagem permite comparar resultados em escalas distintas, desde que a população e o método de cálculo estejam claros.

Os quartis dividem os dados ordenados em quatro partes. $Q_1$ é o primeiro quartil, associado aproximadamente aos $25\%$ inferiores; $Q_2$ é o segundo quartil e coincide com a mediana; $Q_3$ é o terceiro quartil, associado aproximadamente aos $75\%$ inferiores. Os percentis dividem conceitualmente a distribuição em cem partes. Assim, $Q_1=P_{25}$, $Q_2=P_{50}$ e $Q_3=P_{75}$. A palavra ``aproximadamente'' é importante em amostras pequenas, porque posições percentuais podem cair entre observações.

Existem convenções diferentes para calcular quartis e percentis em amostras finitas. Neste capítulo, usaremos um método didático transparente: ordenar; obter a mediana; quando $n$ for par, separar metades de mesmo tamanho; calcular a mediana de cada metade. Quando $n$ for ímpar, excluiremos a mediana geral antes de formar as metades. Softwares podem interpolar de outra maneira e apresentar pequenas diferenças. Isso não é erro automático; o relatório deve declarar a convenção usada.

Considere os dados ordenados $x={2,4,5,7,8,10,12,14}$. Temos $n=8$. A mediana usa o quarto e o quinto valores:
\[
Q_2=\frac{7+8}{2}=7{,}5.
\]
A metade inferior é ${2,4,5,7}$. Sua mediana usa $4$ e $5$:
\[
Q_1=\frac{4+5}{2}=4{,}5.
\]
A metade superior é ${8,10,12,14}$, cuja mediana usa $10$ e $12$:
\[
Q_3=\frac{10+12}{2}=11.
\]

A interpretação precisa acompanhar a conta. $Q_1=4{,}5$ indica que a fronteira do primeiro quarto está entre $4$ e $5$; $Q_2=7{,}5$ separa as metades; $Q_3=11$ separa o quarto superior dos demais. Não significa que exatamente $25\%$ dos valores sejam iguais a $4{,}5$, nem que $4{,}5$ tenha sido observado. Quartil é posição de corte.

Como contraexemplo de convenção, considere $\{1,2,3,4,5\}$, com $n$ ímpar. Pelo método adotado neste capítulo, excluímos a mediana $3$, formamos $\{1,2\}$ e $\{4,5\}$ e obtemos $Q_1=(1+2)/2=1{,}5$ e $Q_3=(4+5)/2=4{,}5$. Outra convenção inclui a mediana nas duas metades: $\{1,2,3\}$ e $\{3,4,5\}$, produzindo $Q_1=2$ e $Q_3=4$. Ambos os pares podem ser corretos sob métodos declarados. O erro seria misturar convenções ou comparar resultados sem registrar como foram obtidos.

Para percentis, adotaremos a posição
\[
L=\frac{k}{100}(n+1),
\]
em que $k$ é o percentil desejado, $n$ é o número de observações e $L$ é a posição teórica na lista ordenada. Se $L$ for inteiro, usamos o valor daquela posição. Se cair entre duas posições, fazemos interpolação linear: percorremos a fração indicada entre os dois valores vizinhos. Essa é uma convenção entre várias, escolhida porque deixa o raciocínio visível.

Calculemos $P_{60}$ no mesmo conjunto com $n=8$. Primeiro,
\[
L=\frac{60}{100}(8+1)=0{,}60\times9=5{,}4.
\]
A posição $5{,}4$ está $0{,}4$ do caminho entre o quinto valor, $8$, e o sexto, $10$. A distância entre eles é $10-8=2$. Tomamos $40\%$ dessa distância: $0{,}4\times2=0{,}8$. Somamos ao valor inferior:
\[
P_{60}=8+0{,}8=8{,}8.
\]
Segundo esta convenção, o corte do percentil 60 é $8{,}8$.

Uma interpretação cuidadosa seria: aproximadamente $60\%$ da distribuição situa-se até o corte $8{,}8$. Não devemos dizer que uma pessoa ``acertou 60\%'' só porque está no percentil 60. Percentual de acertos e percentil são ideias diferentes. Um estudante pode acertar $70\%$ de uma prova e estar no percentil 90 se quase todos acertaram menos; pode acertar $90\%$ e estar no percentil 40 em um grupo de desempenho muito alto.

\begin{GuidedBox}
Aluno, preste atenção aqui: quartis e percentis dependem da ordem e de uma convenção de posição. Ao comparar um cálculo manual com uma ferramenta, confira o método antes de concluir que um dos resultados está errado. Em um relatório, registre a convenção.
\end{GuidedBox}

As posições contam onde os dados estão, mas ainda não dizem quanto eles se espalham. Dois grupos podem compartilhar mediana $7{,}5$ e apresentar comportamentos muito diferentes: ${7,7,7,8,8,8}$ está concentrado; ${1,3,7,8,12,14}$ está espalhado. Precisamos acrescentar medidas de dispersão. Começaremos pela mais simples, a amplitude, e depois construiremos uma medida que utiliza a distância de cada valor ao centro.

\begin{SummaryBox}
Então, aluno, veja só o que você aprendeu aqui: quartis dividem a ordem em quatro regiões e percentis localizam cortes em cem partes conceituais. Eles descrevem posição relativa, exigem dados ordenados e devem vir acompanhados da convenção de cálculo.
\end{SummaryBox}

\section{Amplitude, variância e o motivo do quadrado}

A dispersão tornou-se necessária quando analistas perceberam que o centro, sozinho, não descrevia estabilidade, risco ou diversidade. Duas fábricas podem produzir peças com a mesma medida média e qualidades muito diferentes se uma opera perto do alvo e outra oscila amplamente. Duas turmas podem ter a mesma nota média, mas uma ser homogênea e outra conter resultados muito baixos e muito altos. Medidas de dispersão quantificam essa abertura dos dados.

A medida mais direta é a \textbf{amplitude total}:
\[
A=x_{\max}-x_{\min}.
\]
$x_{\max}$ é o maior valor observado e $x_{\min}$ é o menor. Para os tempos ${18,20,20,22,25,25,30,56}$, temos $x_{\max}=56$ e $x_{\min}=18$. Logo,
\[
A=56-18=38\text{ minutos}.
\]
Do menor ao maior tempo há uma extensão de $38$ minutos.

A amplitude é fácil de explicar, mas utiliza somente dois valores. Se o $56$ fosse substituído por $35$, a amplitude cairia de $38$ para $17$, embora sete observações permanecessem idênticas. Também não diferencia distribuições internas: ${0,0,0,10}$ e ${0,5,5,10}$ possuem amplitude $10$, mas o segundo conjunto ocupa mais o intervalo. Para incorporar todas as observações, calculamos distâncias em relação à média.

Uma primeira tentativa seria fazer a média dos desvios $x_i-\bar{x}$. Porém, já vimos que esses desvios somam zero: os negativos anulam os positivos. Poderíamos usar valores absolutos, e isso produz uma medida válida, o desvio absoluto médio. A tradição da variância escolheu elevar os desvios ao quadrado. O quadrado elimina sinais, dá mais peso a afastamentos grandes e possui propriedades algébricas que sustentam inferência e modelagem. A escolha não é mágica: é uma solução matemática para impedir cancelamento e trabalhar bem com somas.

Quando possuímos toda a população, a \textbf{variância populacional} é
\[
\sigma^2=\frac{1}{N}\sum_{i=1}^{N}(x_i-\mu)^2.
\]
$\sigma^2$ é a variância populacional; $N$ é o tamanho da população; $x_i$ é cada valor; $\mu$ é a média populacional; $x_i-\mu$ é o desvio; e $(x_i-\mu)^2$ é o desvio quadrático. A operação completa é: calcular a média, subtrair a média de cada valor, elevar cada diferença ao quadrado, somar e dividir por $N$.

Considere que as idades ${2,4,4,6}$ sejam toda a população de quatro equipamentos de um laboratório. A média populacional é
\[
\mu=\frac{2+4+4+6}{4}=\frac{16}{4}=4.
\]
Agora abrimos cada desvio e quadrado:
\[
(2-4)^2=(-2)^2=4,
\]
\[
(4-4)^2=0^2=0,
\]
\[
(4-4)^2=0^2=0,
\]
\[
(6-4)^2=2^2=4.
\]
A soma dos quadrados é $4+0+0+4=8$. Como temos a população inteira,
\[
\sigma^2=\frac{8}{4}=2\text{ anos}^2.
\]

\begin{center}
\small
\begin{tabular}{@{}rrr@{}}
\toprule
$x_i$ & $x_i-\mu$ & $(x_i-\mu)^2$ \\
\midrule
2 & $-2$ & 4 \\
4 & 0 & 0 \\
4 & 0 & 0 \\
6 & 2 & 4 \\
\midrule
Soma & 0 & 8 \\
\bottomrule
\end{tabular}
\end{center}

O resultado $2$ está em anos ao quadrado porque elevamos diferenças medidas em anos. Essa unidade dificulta a interpretação direta, mas a variância cumpre uma função central: quantifica o afastamento quadrático médio e penaliza mais os afastamentos grandes. Uma população com variância zero possui todos os valores iguais à média. Quanto maior a variância, maior o espalhamento quadrático na escala considerada.

Dois conjuntos podem ter exatamente a mesma média e dispersões muito diferentes. Considere $A=\{4,5,5,6\}$ e $B=\{1,3,7,9\}$. Ambos somam $20$ e possuem média $5$. Em $A$, os desvios são $-1,0,0,1$, os quadrados somam $2$, $\sigma_A^2=2/4=0{,}5$ e $\sigma_A\approx0{,}71$. Em $B$, os desvios são $-4,-2,2,4$, os quadrados somam $40$, $\sigma_B^2=40/4=10$ e $\sigma_B\approx3{,}16$.

A Figura~\ref{fig:mesma-media-dispersao} mostra os dois conjuntos na mesma escala. A linha vertical marca a média comum. Como a pergunta visual é comparar espalhamento ao redor do mesmo centro, um dotplot preserva cada observação e evita que uma barra agregada esconda a diferença.

\begin{figure}[htbp]
\centering
\begin{tikzpicture}[x=.76cm,y=1.35cm]
\draw[->] (0,.45)--(10.5,.45) node[right] {valor};
\foreach \x in {0,...,10}{\draw[gray!55] (\x,.39)--(\x,.51) node[below] {\x};}
\node[anchor=east] at (0,2) {Conjunto A};
\node[anchor=east] at (0,1) {Conjunto B};
\foreach \x in {4,5,6}{\fill[guidedblueframe] (\x,2) circle (2.8pt);}
\fill[guidedblueframe] (5,2.14) circle (2.8pt);
\foreach \x in {1,3,7,9}{\fill[guidedorangeframe] (\x,1) rectangle +(4pt,4pt);}
\draw[black,thick,dashed] (5,.65)--(5,2.4) node[above] {média $=5$};
\end{tikzpicture}
\caption{Mesma média, mas desvios-padrão muito diferentes. Fonte: elaboração didática.}
\label{fig:mesma-media-dispersao}
\end{figure}

Na Figura~\ref{fig:mesma-media-dispersao}, os pontos de $A$ cercam a média de perto, enquanto os de $B$ ocupam quase toda a escala. Um cartão que mostrasse apenas ``média 5'' faria os processos parecerem equivalentes. A variância e o desvio-padrão respondem justamente ao que o cartão omitiu. Este é o motivo operacional de nunca interpretar centro sem alguma medida de dispersão.

\begin{SummaryBox}
Então, aluno, veja só o que você aprendeu aqui: amplitude usa apenas máximo e mínimo. Variância usa todos os desvios em relação à média, eleva-os ao quadrado para impedir cancelamento e resume a soma desses quadrados. Quando os dados são a população inteira, dividimos por $N$.
\end{SummaryBox}

\section{Variância amostral, o divisor $n-1$ e o desvio-padrão}

Na prática, frequentemente observamos uma amostra e queremos estimar a dispersão da população. Surge uma dificuldade: calculamos os desvios em torno de $\bar{x}$, a média da própria amostra. Essa média foi escolhida pelos mesmos dados e fica exatamente no ponto que minimiza a soma dos quadrados. Consequentemente, os desvios amostrais tendem a parecer um pouco menores do que seriam em torno da média populacional desconhecida. Dividir por $n$ tenderia a subestimar a variância populacional.

A correção tradicional, associada ao desenvolvimento da teoria estatística por Bessel e outros pesquisadores do século XIX, substitui $n$ por $n-1$:
\[
s^2=\frac{1}{n-1}\sum_{i=1}^{n}(x_i-\bar{x})^2.
\]
$s^2$ é a variância amostral; $n$ é o tamanho da amostra; $\bar{x}$ é a média amostral. O termo $n-1$ também expressa um grau de liberdade perdido. Depois que conhecemos $n-1$ desvios e sabemos que sua soma é zero, o último desvio já está determinado; apenas $n-1$ podem variar livremente.

Vamos tratar ${2,4,4,6}$ agora como amostra de uma população maior. A média continua $\bar{x}=4$, e a soma dos desvios quadráticos continua $8$. A diferença está no objetivo. Não estamos descrevendo apenas esses quatro equipamentos; estamos estimando a dispersão da população que eles representam. Logo,
\[
s^2=\frac{8}{4-1}=\frac{8}{3}\approx2{,}67\text{ anos}^2.
\]
Com os mesmos valores, a variância populacional foi $2$ e a amostral é aproximadamente $2{,}67$. Nenhuma é universalmente ``a correta'': a escolha depende de termos população completa ou amostra usada para estimar uma população.

Podemos enxergar o grau de liberdade nos desvios $-2,0,0,2$. Conhecidos os três primeiros e a regra de que a soma deve ser zero, o quarto precisa ser $2$: $-2+0+0+d_4=0$, então $d_4=2$. Não possuímos quatro informações independentes sobre os desvios. Dividir a soma por três compensa, em média, a concentração produzida ao estimar o centro com os próprios dados.

\begin{GuidedBox}
Aluno, preste atenção aqui: não use $n-1$ apenas porque ``a fórmula manda''. Pergunte qual é o papel do conjunto. Se ele contém toda a população definida, descrevemos com $N$. Se é uma amostra usada para estimar a variância populacional, usamos $n-1$.
\end{GuidedBox}

A variância resolve o cancelamento, mas sua unidade ao quadrado é pouco intuitiva. Por isso tomamos a raiz quadrada e obtemos o \textbf{desvio-padrão}. Para a população,
\[
\sigma=\sqrt{\sigma^2};
\]
para a amostra,
\[
s=\sqrt{s^2}.
\]
A raiz desfaz o quadrado da unidade. Se as idades estão em anos, o desvio-padrão volta a anos; se os tempos estão em minutos, volta a minutos. Essa compatibilidade permite comparar diretamente dispersão e média na escala original.

No exemplo populacional, $\sigma^2=2$, então
\[
\sigma=\sqrt{2}\approx1{,}41\text{ anos}.
\]
No exemplo amostral, $s^2=8/3\approx2{,}67$, então
\[
s=\sqrt{2{,}67}\approx1{,}63\text{ anos}.
\]
Podemos dizer que as idades apresentam afastamento quadrático típico de cerca de $1{,}41$ ano em relação à média quando os quatro equipamentos formam a população. A expressão ``típico'' é uma intuição; o desvio-padrão não é a média dos valores absolutos dos desvios.

Façamos um exemplo mais completo com as notas amostrais $x={6,7,7,8,12}$. A soma é $40$ e a média é $40/5=8$. Os desvios são $-2,-1,-1,0,4$. Os quadrados são $4,1,1,0,16$, com soma $22$. A variância amostral é
\[
s^2=\frac{22}{5-1}=\frac{22}{4}=5{,}5.
\]
O desvio-padrão é
\[
s=\sqrt{5{,}5}\approx2{,}35\text{ pontos}.
\]

\begin{center}
\small
\begin{tabular}{@{}rrr@{}}
\toprule
Nota $x_i$ & Desvio $x_i-8$ & Quadrado \\
\midrule
6 & $-2$ & 4 \\
7 & $-1$ & 1 \\
7 & $-1$ & 1 \\
8 & 0 & 0 \\
12 & 4 & 16 \\
\midrule
Soma & 0 & 22 \\
\bottomrule
\end{tabular}
\end{center}

O valor $12$ contribui com $16$ dos $22$ pontos quadráticos, cerca de $73\%$ da soma. Isso mostra a sensibilidade da variância e do desvio-padrão aos extremos. Essa sensibilidade é útil quando grandes desvios representam risco real, como falhas muito demoradas ou custos muito acima do planejado. Quando queremos descrever a região central com maior resistência a extremos, voltamos aos quartis e calculamos o intervalo interquartil.

\begin{SummaryBox}
Então, aluno, veja só o que você aprendeu aqui: $n-1$ corrige a tendência de a amostra parecer concentrada em torno da média que ela própria estimou. O desvio-padrão é a raiz da variância e devolve a medida à unidade original, mas continua sensível a afastamentos grandes.
\end{SummaryBox}

\section{Coeficiente de variação: comparar dispersão relativa com cuidado}

O desvio-padrão conserva a unidade original, o que é excelente para interpretar uma variável, mas dificulta certas comparações. Um desvio de R\$ 500 pode ser pequeno diante de salário médio de R\$ 10.000 e grande diante de renda média de R\$ 800. Para responder ``quanto a dispersão representa em relação ao centro?'', usamos o \textbf{coeficiente de variação}, difundido no desenvolvimento da biometria e da comparação de séries com escalas distintas. Ele transforma o desvio-padrão em uma razão relativa à média.

Para uma amostra com média positiva, o coeficiente de variação é
\[
CV=\frac{s}{\bar{x}}\times100\%.
\]
$s$ é o desvio-padrão amostral, $\bar{x}$ é a média amostral e a multiplicação por $100\%$ expressa a razão em percentagem. Para uma população, podemos escrever $CV=(\sigma/\mu)\times100\%$. Como numerador e denominador possuem a mesma unidade, ela se cancela. O CV não informa minutos ou reais; informa dispersão relativa ao tamanho do centro.

Considere dois processos. O processo A possui média $50$ ms e desvio-padrão $5$ ms. O processo B possui média $10$ ms e desvio-padrão $2$ ms. Em termos absolutos, A tem maior desvio: $5>2$. Porém,
\[
CV_A=\frac{5}{50}\times100\%=10\%,
\]
\[
CV_B=\frac{2}{10}\times100\%=20\%.
\]
O processo B varia menos milissegundos, mas varia o dobro em relação à própria média. A pergunta ``qual oscila mais?'' precisa dizer se interessa variação absoluta ou relativa.

A razão só é interpretável quando o zero da escala significa ausência da quantidade e a média está suficientemente distante de zero. Massa, tempo, distância e receita positiva costumam admitir comparação relativa. Temperatura em graus Celsius é um contraexemplo: $20^\circ$C não representa o dobro de calor de $10^\circ$C, porque o zero Celsius é convencional. Converter para Fahrenheit mudaria o CV embora o fenômeno fosse o mesmo. Nesses casos, DP e outras descrições são preferíveis.

Médias próximas de zero também tornam o CV instável. Se um retorno médio é $0{,}2\%$ com desvio $2\%$, o CV seria $1.000\%$; se a média muda ligeiramente para $-0{,}1\%$, o sinal muda e a razão perde a leitura usual. Para dados que podem ser negativos, como lucro líquido ou variação de temperatura, o CV pode não responder a uma pergunta coerente. Não basta a fórmula aceitar a divisão: a escala precisa sustentar a interpretação.

O CV também não substitui a visualização nem resolve assimetria. Dois processos podem ter CV semelhante e distribuições com caudas diferentes. Ao publicá-lo, informe média, desvio-padrão, unidade e tamanho da amostra. Assim, o leitor consegue reconstruir a razão e avaliar se o centro positivo e a escala de razão tornam a comparação válida.

Considere agora dois times que possuem tempo médio próximo de nove minutos e precisam cumprir um SLA de até doze minutos na maior parte dos atendimentos. A tabela reúne centro, dispersão e tamanho. O time A parece mais previsível pelo DP e pelo IQR; o time B possui metade central e afastamento quadrático mais largos. Entretanto, $n=8$ em B torna sua estimativa mais sensível a poucos casos. A decisão não é premiar A automaticamente, mas investigar proporção acima do SLA, composição e período comparável.

\begin{center}
\small
\begin{tabular}{@{}lrrrr@{}}
\toprule
\textbf{Time} & $n$ & \textbf{Média} & \textbf{DP} & \textbf{IQR} \\
\midrule
A & 40 & 9,1 min & 1,2 min & 1,5 min \\
B & 8 & 9,0 min & 3,8 min & 4,0 min \\
\bottomrule
\end{tabular}
\end{center}

Comparações por grupo devem publicar $n$, período e unidade de análise. Uma média baseada em duas pessoas e outra baseada em duzentas podem ter o mesmo valor, mas não a mesma estabilidade. Também precisamos confirmar que os grupos enfrentaram tipos de chamado comparáveis. A regra editorial é simples: nunca apresente média ou mediana por grupo sem informar quantas observações a sustentam e qual decisão será tomada.

DP e IQR respondem a sensibilidades diferentes. Quando a distribuição é aproximadamente simétrica e sem extremos dominantes, média e DP preservam bem centro e espalhamento. Quando há assimetria ou casos afastados legítimos, mediana e IQR descrevem melhor a região central resistente. Em distribuição bimodal, nenhum par encerra a análise: investigamos subgrupos. O IQR não mede simetria; apenas a largura da metade central.

\begin{GuidedBox}
Aluno, preste atenção aqui: use CV para comparar dispersão relativa quando a média positiva e o zero significativo permitem a razão. Não use por hábito em Celsius, escalas arbitrárias, médias próximas de zero ou valores que atravessam zero.
\end{GuidedBox}

O coeficiente de variação encerra a família baseada em média e desvio-padrão. Para uma leitura resistente da metade central e para uma representação que sinalize observações afastadas, voltaremos aos quartis. O IQR e o boxplot completarão a comparação entre dispersão sensível e dispersão robusta.

\begin{SummaryBox}
Então, aluno, veja só o que você aprendeu aqui: DP mede dispersão na unidade; CV divide essa dispersão pela média para comparações relativas. O ganho de comparabilidade vem acompanhado de uma condição: a razão precisa ter significado no fenômeno.
\end{SummaryBox}

\section{IQR e boxplot: posição e dispersão em uma única leitura}

O intervalo interquartil, conhecido pela sigla inglesa IQR, surgiu como uma forma robusta de descrever a largura da metade central dos dados. Sua fórmula é
\[
IQR=Q_3-Q_1.
\]
Como utiliza os quartis e ignora diretamente os $25\%$ inferiores e os $25\%$ superiores ao medir a largura, sofre menos influência de valores extremos do que a amplitude e o desvio-padrão. Ele não substitui todas as medidas; responde especificamente quanto espaço a metade central ocupa.

Retomemos os tempos ${18,20,20,22,25,25,30,56}$. A mediana é $(22+25)/2=23{,}5$. A metade inferior é ${18,20,20,22}$, então
\[
Q_1=\frac{20+20}{2}=20.
\]
A metade superior é ${25,25,30,56}$, então
\[
Q_3=\frac{25+30}{2}=27{,}5.
\]
Logo,
\[
IQR=27{,}5-20=7{,}5\text{ minutos}.
\]
A metade central dos tempos ocupa uma faixa de $7{,}5$ minutos.

Compare IQR e amplitude. A amplitude é $56-18=38$ minutos, muito influenciada pelo maior valor. O IQR é $7{,}5$, refletindo a concentração do miolo entre $20$ e $27{,}5$. Juntas, as medidas sugerem que existe uma cauda superior ou observação afastada. Não devemos apagar o $56$ apenas porque é distante. Primeiro investigamos se ele é válido e o que ocorreu: dificuldade técnica, interrupção, acessibilidade, erro de registro ou variação legítima.

O boxplot foi popularizado no século XX pelo estatístico John Tukey como parte de uma abordagem exploratória que valorizava olhar para padrões antes de impor modelos. Ele resume cinco posições: mínimo não atípico, $Q_1$, mediana, $Q_3$ e máximo não atípico. A caixa vai de $Q_1$ a $Q_3$; uma linha marca a mediana; hastes alcançam os valores observados mais extremos ainda compatíveis com o critério adotado; observações além dele são desenhadas separadamente.

O critério usual calcula cercas de referência:
\[
LI=Q_1-1{,}5\times IQR,
\qquad
LS=Q_3+1{,}5\times IQR.
\]
$LI$ é a cerca inferior e $LS$ a cerca superior. Para os tempos,
\[
LI=20-1{,}5\times7{,}5=20-11{,}25=8{,}75,
\]
\[
LS=27{,}5+1{,}5\times7{,}5=27{,}5+11{,}25=38{,}75.
\]
Qualquer valor abaixo de $8{,}75$ ou acima de $38{,}75$ será sinalizado pelo critério de Tukey. O $56$ está acima da cerca superior.

As cercas não são necessariamente as extremidades desenhadas das hastes. A haste inferior alcança o menor valor observado que permanece dentro das cercas: $18$. A haste superior alcança o maior valor observado dentro das cercas: $30$. O $56$ aparece como ponto separado. O resumo do boxplot é, portanto: haste inferior $18$, $Q_1=20$, mediana $23{,}5$, $Q_3=27{,}5$, haste superior $30$ e um ponto sinalizado em $56$.

\begin{center}
\small
\begin{tabular}{@{}lr@{}}
\toprule
\textbf{Elemento} & \textbf{Valor} \\
\midrule
Haste inferior & 18 \\
$Q_1$ & 20 \\
Mediana & 23,5 \\
$Q_3$ & 27,5 \\
Haste superior & 30 \\
Valor sinalizado & 56 \\
\bottomrule
\end{tabular}
\end{center}

Um ponto fora das cercas não é automaticamente erro e não deve ser chamado de impossível. O termo \textit{outlier} descreve uma observação afastada segundo um critério, não uma sentença sobre sua origem. Se $56$ foi digitado no lugar de $26$, corrigimos com evidência na fonte. Se representa um estudante que enfrentou barreira de acessibilidade, removê-lo esconderia um problema decisivo. Se é uma ocorrência legítima rara, podemos apresentar estatísticas com e sem o caso, explicando a razão, mas não apagá-lo silenciosamente.

O formato do boxplot também informa assimetria. Neste exemplo, a distância de $Q_3=27{,}5$ até a haste superior $30$ é pequena, mas existe um ponto muito distante em $56$. A média $27$ fica acima da mediana $23{,}5$, coerente com a influência da cauda superior. Um boxplot não mostra bem multimodalidade nem todos os detalhes da distribuição; conjuntos diferentes podem produzir caixas semelhantes. Por isso ele deve ser acompanhado de contexto, contagem e, quando necessário, outra representação.

A Figura~\ref{fig:boxplot-efeito-extremo} compara a versão sem o $56$, $\{18,20,20,22,25,25,30\}$, com a versão completa. No primeiro conjunto, pela convenção declarada de excluir a mediana nas metades ímpares, temos $Q_1=20$, mediana $22$ e $Q_3=25$; as hastes vão de $18$ a $30$. No segundo, a caixa vai de $20$ a $27{,}5$, a mediana é $23{,}5$, as hastes ainda terminam em $18$ e $30$ e o $56$ é desenhado separadamente.

\begin{figure}[htbp]
\centering
\begin{tikzpicture}[x=.125cm,y=1.25cm]
\draw[->] (5,0) -- (61,0) node[right] {minutos};
\foreach \x in {10,20,30,40,50,60}{\draw (\x,.08)--(\x,-.08) node[below] {\x};}
\node[anchor=east] at (5,1) {sem 56};
\node[anchor=east] at (5,2.2) {com 56};
\draw[thick] (18,1)--(20,1) (25,1)--(30,1);
\draw[thick,fill=guidedblue] (20,.72) rectangle (25,1.28);
\draw[very thick,guidedblueframe] (22,.72)--(22,1.28);
\draw[thick] (18,.82)--(18,1.18) (30,.82)--(30,1.18);
\draw[thick] (18,2.2)--(20,2.2) (27.5,2.2)--(30,2.2);
\draw[thick,fill=guidedorange] (20,1.92) rectangle (27.5,2.48);
\draw[very thick,guidedorangeframe] (23.5,1.92)--(23.5,2.48);
\draw[thick] (18,2.02)--(18,2.38) (30,2.02)--(30,2.38);
\filldraw[guidedorangeframe] (56,2.2) circle (2.2pt) node[above] {$56$};
\end{tikzpicture}
\caption{O valor 56 altera quartis e mediana moderadamente, mas aparece separado pelo critério do IQR. Fonte: elaboração didática.}
\label{fig:boxplot-efeito-extremo}
\end{figure}

Na Figura~\ref{fig:boxplot-efeito-extremo}, a inclusão do extremo não estica a haste até $56$, porque a haste termina no maior valor observado dentro da cerca. Também vemos que o IQR muda de $25-20=5$ para $27{,}5-20=7{,}5$: resistência não significa imobilidade. Quartis reagem menos que média e amplitude, mas ainda refletem alterações de posição. Essa leitura impede dois erros opostos: imaginar que o boxplot ignora extremos ou imaginar que todo ponto separado deve ser removido.

\begin{GuidedBox}
Aluno, preste atenção aqui: as cercas de $1{,}5\times IQR$ são um sinalizador exploratório, não uma regra automática de exclusão. O próximo passo é investigar a observação e sua origem. Estatística ajuda a formular a pergunta; o contexto sustenta a decisão.
\end{GuidedBox}

O boxplot completa a ponte do capítulo. Frequências organizaram ocorrências; média, mediana e moda resumiram centros diferentes; quartis e percentis localizaram posições; amplitude, variância, desvio-padrão e IQR mediram dispersões distintas. Agora dispomos de uma pequena gramática para descrever dados antes de produzir qualquer visual mais elaborado ou inferência.

\begin{SummaryBox}
Então, aluno, veja só o que você aprendeu aqui: o IQR mede a largura da metade central. O boxplot reúne quartis, mediana, hastes e pontos sinalizados. Ele ajuda a enxergar posição, dispersão e assimetria, mas não substitui a investigação do contexto.
\end{SummaryBox}

\section{Caso integrado e síntese da aula de três horas}

Vamos fechar reconstruindo a análise de um novo conjunto sem atalhos. Uma equipe registrou, para oito atendentes, $4$, $5$, $5$, $6$, $6$, $7$, $7$ e $20$ chamados resolvidos no turno. Suponha que esses oito formem toda a equipe daquele turno; portanto, tratamos o conjunto como população para descrevê-lo. A unidade de observação é o atendente e a variável ``chamados resolvidos'' é quantitativa discreta. Antes de calcular, confirmamos que cada atendente aparece uma vez e que todos trabalharam o mesmo período.

As frequências absolutas são: $4$ aparece uma vez; $5$, duas; $6$, duas; $7$, duas; $20$, uma. As frequências relativas são $1/8=12{,}5\%$ para $4$ e $20$, e $2/8=25\%$ para $5$, $6$ e $7$. A soma é $12{,}5+25+25+25+12{,}5=100\%$. A distribuição é multimodal, com modas $5$, $6$ e $7$. Essa repetição mostra um grupo central, enquanto $20$ merece investigação.

A soma dos valores é $4+5+5+6+6+7+7+20=60$. A média populacional é
\[
\mu=\frac{60}{8}=7{,}5\text{ chamados}.
\]
Como os dados já estão ordenados e $n=8$, a mediana usa a quarta e a quinta posições, ambas iguais a $6$:
\[
Q_2=\frac{6+6}{2}=6.
\]
A média $7{,}5$ fica acima da mediana $6$ por causa do valor $20$. Dizer que o atendente ``típico'' resolveu $7{,}5$ chamados sem informar a assimetria seria incompleto.

Para os quartis, a metade inferior é ${4,5,5,6}$, então $Q_1=(5+5)/2=5$. A metade superior é ${6,7,7,20}$, então $Q_3=(7+7)/2=7$. O IQR é
\[
IQR=7-5=2.
\]
As cercas são $LI=5-1{,}5\times2=2$ e $LS=7+1{,}5\times2=10$. Assim, $20$ é sinalizado. As hastes alcançam $4$ e $7$; a caixa vai de $5$ a $7$ e a mediana fica em $6$.

Agora calculamos a variância populacional. Em relação a $\mu=7{,}5$, os desvios são $-3{,}5$, $-2{,}5$, $-2{,}5$, $-1{,}5$, $-1{,}5$, $-0{,}5$, $-0{,}5$ e $12{,}5$. Seus quadrados são $12{,}25$, $6{,}25$, $6{,}25$, $2{,}25$, $2{,}25$, $0{,}25$, $0{,}25$ e $156{,}25$. A soma é $186$. Logo,
\[
\sigma^2=\frac{186}{8}=23{,}25\text{ chamados}^2,
\qquad
\sigma=\sqrt{23{,}25}\approx4{,}82\text{ chamados}.
\]
O $20$ domina a soma quadrática, assim como já suspeitávamos pelo boxplot.

Se os mesmos oito atendentes fossem uma amostra de uma rede maior, manteríamos média e desvios, mas usaríamos
\[
s^2=\frac{186}{8-1}=\frac{186}{7}\approx26{,}57,
\qquad
s\approx5{,}15.
\]
A troca não ocorre porque o conjunto mudou, mas porque mudou a pergunta: descrição de toda a equipe do turno ou estimação de uma população mais ampla. Esse é o tipo de decisão que uma ferramenta não pode adivinhar a partir dos números.

A leitura integrada é mais rica do que qualquer medida isolada. A moda aponta concentração em $5$, $6$ e $7$; a mediana $6$ representa a posição central; a média $7{,}5$ revela a influência do total produzido pelo caso $20$; o IQR $2$ descreve um miolo estreito; o desvio-padrão $4{,}82$ reage ao afastamento grande; o boxplot sinaliza o caso para investigação. Se o atendente de $20$ recebeu chamados muito simples ou teve registro duplicado, a comparação individual não é justa. Se ele criou uma prática eficiente, o caso pode gerar aprendizado.

Os primeiros noventa minutos do encontro sustentam população, amostra, variáveis, frequências e médias com contas reconstruídas. A transição até a chamada trabalha mediana, moda e comparação de centros. Depois das 20h30, quartis, percentis, amplitude, variância, desvio-padrão, IQR e boxplot completam a análise. O fechamento integrado existe para que cada símbolo retorne à pergunta original; a fórmula não deve permanecer separada da decisão que a tornou necessária.

Na próxima semana, ampliaremos a estatística básica para probabilidade, distribuições e relação entre variáveis. Média e variância reaparecerão como parâmetros de distribuições; quartis ajudarão a ler formas; e a lógica dos desvios preparará covariância e correlação. Não começaremos outra coleção de fórmulas: continuaremos a mesma história de como dados variam e de como essa variação pode ser descrita com responsabilidade.

\begin{SummaryBox}
Então, aluno, veja só o que você construiu: partiu do significado da linha, organizou frequências, calculou centros, localizou quartis, mediu dispersão e reuniu tudo no boxplot. A boa análise não escolhe um número favorito; compara medidas, explica diferenças e retorna ao contexto antes de decidir.
\end{SummaryBox}
